%==========================================================================%
%  FrankWolfeSolver -- Technical report / design note
%
%  Describes the porting into SMS++ of the Frank-Wolfe (conditional gradient)
%  algorithms, the design choices, the novel elements, and in particular the
%  "qualifying feature" of the two problems (intCvxComb).
%
%  Compile with pdflatex (two passes for the cross-references).
%==========================================================================%
\documentclass[11pt,a4paper]{article}

\usepackage[T1]{fontenc}
\usepackage[utf8]{inputenc}
\usepackage[english]{babel}
\usepackage[margin=2.6cm]{geometry}
\usepackage{amsmath,amssymb,amsthm}
\usepackage{booktabs}
\usepackage{xcolor}
\usepackage{listings}
\usepackage{enumitem}
\usepackage[colorlinks=true,linkcolor=blue!55!black,citecolor=blue!55!black,
            urlcolor=blue!55!black]{hyperref}

%--------------------------------------------------------------------------%
% notation
%--------------------------------------------------------------------------%
\newcommand{\R}{\mathbb{R}}
\newcommand{\conv}{\operatorname{conv}}
\newcommand{\dom}{\operatorname{dom}}
\newcommand{\clamp}{\operatorname{clamp}}
\newcommand{\argmin}{\operatorname*{arg\,min}}
\newcommand{\ip}[2]{\langle #1,\, #2\rangle}
\newcommand{\ffat}{f_{\mathrm{father}}}
\newcommand{\code}[1]{\texttt{#1}}
\newcommand{\Pone}{(\mathrm{P}_1)}
\newcommand{\Ptwo}{(\mathrm{P}_2)}

\theoremstyle{plain}
\newtheorem{proposition}{Proposition}
\theoremstyle{remark}
\newtheorem*{remark}{Remark}

% soak up the few overfull boxes from long ttfamily identifiers
\emergencystretch=3em
\hbadness=4000

\lstdefinestyle{cfg}{basicstyle=\ttfamily\small,frame=single,
  framesep=4pt,rulecolor=\color{black!30},backgroundcolor=\color{black!3},
  columns=fullflexible,keepspaces=true,xleftmargin=2pt}

\title{\bfseries \code{FrankWolfeSolver}:\\
       a Frank-Wolfe \code{CDASolver} for SMS++,\\
       and the \emph{two problems} of a single decomposition}
\author{Technical documentation of the \code{FrankWolfeSolver} module}
\date{\today}

\begin{document}
\maketitle

\begin{abstract}
\noindent
\code{FrankWolfeSolver} is an SMS++ \code{CDASolver} that ports into C++ the
Frank-Wolfe (conditional gradient) family of algorithms of the Julia package
\href{https://github.com/ZIB-IOL/FrankWolfe.jl}{FrankWolfe.jl}, adapting them to
the SMS++ \code{Block} tree: it minimizes a composite objective over a ``father''
Block whose feasible region is the \emph{product of the convex hulls} of the
sub-Block regions, using one \code{:Solver} per sub-Block as a \emph{Linear
Minimization Oracle} (LMO). This document describes the work done, the design
choices, and the novel elements. The main conceptual contribution is the
observation --- and the implementation --- that, \emph{with the same
decomposition and the same oracle}, a single algorithmic parameter
(\code{intCvxComb}) makes the solver solve \emph{two distinct problems}: the
``true'' composite objective over the convex hull, or its (stronger)
Dantzig-Wolfe / perspective-cut relaxation. We discuss the history and motivation
of this distinction, prove the relationship between them (Jensen's inequality,
convex envelope), and detail \emph{exactly what changes in the algorithm} to
solve one or the other. We close with the lazy handling of \code{Modification}
with warm-start, and with the experimental validation.
\end{abstract}

%==========================================================================%
\section{Introduction and context}\label{sec:intro}
%==========================================================================%

\code{FrankWolfeSolver} registers to a ``father'' \code{Block} with the following
structure:
\begin{itemize}[nosep]
  \item an \code{FRealObjective} whose \code{Function} is a \code{C05Function}
        (the \emph{linking} function, of which the \emph{diagonal linearization},
        i.e.\ the gradient, is used);
  \item an arbitrary number of \emph{sub-Block}, holding \emph{all} the variables
        and constraints (the father has none of its own);
  \item each sub-Block has an \code{FRealObjective} containing a
        \code{LinearFunction}, \code{DQuadFunction} or \code{QuadFunction}.
\end{itemize}
To each sub-Block \code{FrankWolfeSolver} registers a \code{:Solver}, which is its
LMO; the solver names come from the \code{Configuration} parameters, exactly as in
\code{LagrangianDualSolver} (LDS). The overall feasible region is the
\emph{product} of the children's regions,
\begin{equation}\label{eq:X}
  X \;=\; \conv(X_1)\times\conv(X_2)\times\cdots\times\conv(X_k),
\end{equation}
where $X_j$ is the region of sub-Block $j$. Any integrality constraints in the
children are \emph{not} looked at by the solver: combining integer solutions as
convex combinations effectively minimizes over the convex hull $\conv(X_j)$.

\paragraph{Analogy with \code{LagrangianDualSolver}.}
The kinship is direct and drives much of the \emph{plumbing}, copied
\emph{verbatim} from LDS where possible (registration/deregistration of the
children's solvers, \code{Configuration} cache, per-child \code{BlockSolverConfig}
parameters, snapshot/restore of the children's objective, \code{Modification}
handling):

\begin{center}
\begin{tabular}{@{}lll@{}}
\toprule
Aspect & \code{LagrangianDualSolver} & \code{FrankWolfeSolver}\\
\midrule
Linking in the father & linear \emph{constraints}      & \emph{objective} \code{C05Function}\\
What it does with it  & dualizes (multipliers)          & linearizes (gradient)\\
Solver for the children & one per sub-Block             & one per sub-Block ($=$ LMO)\\
\bottomrule
\end{tabular}
\end{center}

\paragraph{Status.} \emph{v1} implements ``vanilla'' Frank-Wolfe; the active-set
variants (Away-step and Blended Pairwise) are implemented and tested; the
\code{Modification} handling with warm-start and the infeasibility propagation are
complete.

%==========================================================================%
\section{The composite objective and the three oracle semantics}\label{sec:composite}
%==========================================================================%

The three possible semantics for how the children's objective enters the problem
are special cases of a single \emph{generalized Frank-Wolfe / composite
conditional gradient} scheme, in which the oracle (i.e.\ the child's solver)
handles \emph{exactly} a separable part of the objective. The total minimized
objective is
\begin{equation}\label{eq:F}
  F(x) \;=\; \ffat(x) \;+\; \sum_{j} h_j(x_j), \qquad
  h_j(y) \;=\; \alpha\,\ip{c_j}{y} + \beta\,q_j(y),
\end{equation}
where $c_j,q_j$ are the \emph{original} linear and quadratic parts of child $j$'s
objective, $g=\nabla \ffat(x)$ and $g_j$ is its restriction to child $j$. The
smooth part $\ffat$ is linearized; the $h_j$ stay exact in the oracle. The three
options are the values of $(\alpha,\beta)$, selected by the integer parameter
\code{intLMOObj}:

\begin{center}
\begin{tabular}{@{}clccl@{}}
\toprule
\code{intLMOObj} & name & $(\alpha,\beta)$ & child $j$'s oracle solves & LMO type\\
\midrule
0 & \code{LMOLinear} (default) & $(0,0)$ & $\min\ \ip{g_j}{v_j}$ & LP/MILP\\
1 & \code{LMOQuad}             & $(0,1)$ & $\min\ \ip{g_j}{v_j} + q_j(v_j)$ & QP/MIQP\\
2 & \code{LMOFull}             & $(1,1)$ & $\min\ \ip{c_j+g_j}{v_j} + q_j(v_j)$ & QP/MIQP\\
\bottomrule
\end{tabular}
\end{center}

\paragraph{Block-tree semantics.} The total objective of a Block tree in SMS++ is
the \emph{sum} of all the objectives in the tree: this is exactly what a
\code{:MILPSolver} on the father minimizes (father objective $+\sum_j$ children's
objectives, recursively). Therefore \textbf{\code{LMOFull} is the mode that
coincides with the ``true'' Block-tree objective} and with what the monolithic
\code{:MILPSolver} computes. \code{LMOLinear} (children ignored) coincides with
the monolithic one only if the children have a null objective; \code{LMOQuad}
only if they have a purely quadratic one.

\paragraph{Mechanics: snapshot $+$ scatter $+$ restore.}
In the style of \code{LagBFunction::cleanup\_inner\_objective} of LDS: in
\code{set\_Block} each child's original objective is \emph{snapshot}ted ($c_j$,
and $q_j$ if present); at every iteration the linear coefficients of the child's
objective are \emph{scattered} to $\alpha\,c_j+g_j$ via
\code{modify\_linear\_coefficients} (the quadratic part is never touched: it stays
if $\beta=1$, it is zeroed once and for all if $\beta=0$); at the end of the run
the original objective is \emph{restored}, so the Block stays clean. Scattering
onto a \emph{subset} of variables (the subset path) is essential when the father
touches only some of the child's variables --- the case of the DP formulation
with perspective cuts, where the cut rows must not be touched.

\paragraph{The gap, identical in all modes.}
The generalized Frank-Wolfe gap collapses into a single form:
\begin{equation}\label{eq:gap}
  \mathrm{gap}(x) \;=\; \sum_j \big[\, M_j(x_j) - M_j(v_j) \,\big] \;\geq\; F(x)-F^*,
\end{equation}
where $M_j$ is the child's objective \emph{as modified} ($\ip{\alpha c_j+g_j}{\cdot}
+\beta q_j$): $M_j(v_j)$ is the optimal value returned by the child's solver and
$M_j(x_j)$ is the same objective evaluated at the current iterate. In
\code{LMOLinear} it reduces to $\ip{g}{x}-\ip{g}{v}$. A single routine for the
stopping criterion, independent of the mode.

%==========================================================================%
\section{The two problems: the heart of the contribution}\label{sec:twoproblems}
%==========================================================================%

\subsection{History and motivation}

The observation arises from wanting to use \code{FrankWolfeSolver} as a
\emph{decomposition alternative} to an explicit Dantzig-Wolfe (DW) or
perspective-cut (P/C) reformulation. In the applications of interest (e.g.\ unit
commitment with \code{ThermalUnitBlock}) the sub-Block have a \emph{convex
nonlinear} objective (typically quadratic) and a combinatorial domain; the bound
one wants to compute is that of the \emph{continuous relaxation} obtained by
describing the convex hull of the integer solutions --- exactly what a DP
formulation with perspective cuts computes. The natural question was: \emph{does
Frank-Wolfe on this father compute that bound, or a different one?}

The answer is that \emph{it depends on how the children's term $h_j$ is
accounted for}, and that the two accountings correspond to two mathematically
distinct problems, both legitimate and useful. The same oracle serves both; which
one is solved is a choice --- and it has become a \emph{qualifying feature} of the
solver, governed by the \code{intCvxComb} parameter.

\subsection{The two definitions}

Over the same region $X$ of \eqref{eq:X}, with the vertices $v_k$ generated by the
oracle and the convex weights $\lambda$:
\begin{align}
  \Pone\quad & \min_{x\in X}\ \ffat(x) + \sum_j h_j(x_j),
    \label{eq:P1}\\[2pt]
  \Ptwo\quad & \min_{\lambda}\ \ffat\!\Big(\textstyle\sum_k \lambda_k v_k\Big)
    + \sum_j \sum_k \lambda_k^{\,j}\, h_j\big(v_k^{\,j}\big).
    \label{eq:P2}
\end{align}
$\Pone$ evaluates each child cost $h_j$ \emph{at the} (generally fractional)
\emph{iterate} $x_j$; $\Ptwo$ accounts for it as the \emph{convex combination of
the vertex costs}, i.e.\ it replaces $h_j$ by its \emph{convex envelope} over the
generated vertices.

\subsection{The relationship between the two}

Let $f_1,f_2$ be the two objectives as functions of $x\in X$.

\begin{proposition}\label{prop:jensen}
$f_1$ and $f_2$ agree on the vertices of $X$. By convexity of $h_j$ (Jensen's
inequality), $f_2(x)\geq f_1(x)$ inside $X$, and $f_2$ is the \emph{convex
envelope} of $f_1$. Consequently $\mathrm{val}\,\Ptwo \geq \mathrm{val}\,\Pone$:
$\Ptwo$ is the \emph{stronger} relaxation. The two coincide if and only if all the
$h_j$ are linear.
\end{proposition}

\begin{proof}[Proof (sketch)]
On a vertex $\lambda$ is degenerate, so the two children terms coincide. For
$x_j=\sum_k\lambda_k^j v_k^j$ with $h_j$ convex, Jensen gives
$h_j(x_j)\le\sum_k\lambda_k^j h_j(v_k^j)$, whence $f_1\le f_2$ pointwise; the
difference is zero exactly when each $h_j$ is affine on the active vertices, i.e.\
(for every choice of vertices) when the $h_j$ are linear.
\end{proof}

\paragraph{Dantzig-Wolfe / perspective-cut reading.}
$\Ptwo$ is precisely the \emph{Dantzig-Wolfe / simplicial-decomposition} bound:
Frank-Wolfe generates the columns (the sub-Block vertices) and the ``master''
recombines them with their \emph{true} costs. When $\conv(X_j)$ is the convex hull
of the integer solutions of the sub-Block and $h_j$ is a convex-quadratic cost,
$\Ptwo$ \emph{is} the value of the perspective reformulation / of the
perspective-cut (P/C) bound of that sub-Block. In other words:
\textbf{\code{FrankWolfeSolver} in mode $\Ptwo$ is a decomposition alternative to
an explicit DW / P-C reformulation}, without having to write the master nor the
columns by hand.

\subsection{What changes in the algorithm}\label{sec:whatchanges}

The crucial point --- and what makes the feature ``free'' --- is that
\textbf{the oracle (LMO) is identical in the two modes}: it always returns
vertices that account for the exact $h_j$. Only the value/gap/line-search
\emph{bookkeeping} changes. In detail:

\begin{enumerate}[leftmargin=1.4em]
\item \textbf{Per-atom stored cost.} Each atom of the active set stores
  $c_i = \sum_j h_j(\text{atom}_j)$, the child cost \emph{at the vertex}. The
  children term of the current value is then
  \[
    \underbrace{\bar c = \sum_i \lambda_i\, c_i}_{\Ptwo\ (\code{eObjCvxComb})}
    \qquad\text{versus}\qquad
    \underbrace{\sum_j h_j(x_j)}_{\Pone\ (\code{eObjAtX})},
  \]
  and the current value is
  $\mathrm{cost\_x} = (\text{the}\ \ip{g}{x}\ \text{piece}) + \big(\text{cvx}\,?\,
  \bar c:\sum_j h_j(x_j)\big)$. For an \emph{aggregated} atom (merge of atoms when
  the active set exceeds \code{intMaxAtoms}), $c_i$ is the convex combination of
  the merged atoms' costs: it is the correct quantity for $\Ptwo$, already cached.

\item \textbf{Line search.} Along the segment $x+\gamma(v-x)$:
  \begin{itemize}[nosep]
    \item in $\Ptwo$ the children term is \emph{linear in $\gamma$} (it
      interpolates $c_i$ between the two atoms): the exact line search reduces to
      the father-only quadratic case, in closed form,
      \begin{equation}\label{eq:ls}
        \gamma^\star = \clamp\!\Big(-\tfrac{\ip{\nabla F}{d}}{2Q},\,0,\,1\Big),\quad
        d=v-x,\ \ Q=\tfrac12\ip{d}{A d},
      \end{equation}
      with $A$ the (static) father Hessian, \emph{cached} in \code{set\_Block};
    \item in $\Pone$ the children term is $h_j\big(x_j(\gamma)\big)$,
      \emph{quadratic (nonlinear) in $\gamma$} when $h_j$ is quadratic.
  \end{itemize}

\item \textbf{Different iterates, not just a different value.} Since the iterates
  are driven by the line search, the two modes follow \emph{different
  trajectories} and converge to the two different optima $\mathrm{val}\,\Pone$ and
  $\mathrm{val}\,\Ptwo$: it is not a mere change of the reported value.

\item \textbf{Gap and stopping criterion} use \texttt{cost\_x} consistent with the
  mode; in $\Ptwo$ the exact-line-search ``gate'' is relaxed, because the model
  along the segment is linear-in-the-weights.
\end{enumerate}

The parameter selecting the mode is:
\begin{center}
\begin{tabular}{@{}clll@{}}
\toprule
\code{intCvxComb} & enum & solves & children term of the value\\
\midrule
1 & \code{eObjCvxComb} \emph{(default)} & $\Ptwo$ (DW / P-C) & $\bar c=\sum_i\lambda_i c_i$ (linear in $\gamma$)\\
0 & \code{eObjAtX}                       & $\Pone$ (composite) & $\sum_j h_j(x_j)$ (nonlinear in $\gamma$)\\
\bottomrule
\end{tabular}
\end{center}

\begin{remark}
Default \code{eObjCvxComb}: it is the stronger bound and the one matching the
P/C, the primary use. v1 caveat: the \emph{exact} line search of $\Pone$ with a
\emph{nonlinear} $h_j$ is not yet implemented (it would require solving a
quadratic in the children term along the segment); in that case $\Pone$ falls back
on the agnostic $2/(t+2)$ rule. For linear $h_j$ the two modes coincide
(Proposition~\ref{prop:jensen}) and the issue does not arise.
\end{remark}

%==========================================================================%
\section{The algorithmic scheme}\label{sec:algo}
%==========================================================================%

\paragraph{Variants (\code{intAlgorithm}).}
\code{AlgVanilla} (0, no active set), \code{AlgAwayStep} (1, Away-step
Frank-Wolfe) and \code{AlgBPCG} (2, Blended Pairwise Conditional Gradient). The
two active-set variants maintain the set of \emph{atoms} (generated vertices) and
their weights, and are bounded by \code{intMaxAtoms} via aggregation of the least
active ones.

\paragraph{LMO and parallelism.}
Each iteration: the father gradient is extracted, \emph{scattered} onto the
children's objectives (Sec.~\ref{sec:composite}), each child's \code{:Solver} is
run (\code{run\_LMOs}) and the vertex $v_j$ is materialized via
\code{get\_var\_solution}. The LMOs are independent across children:
\code{intMaxThread} enables a persistent bulk-synchronous thread pool.

\paragraph{Line search (\code{intLineSearch}).}
\code{LSAuto} (exact if the father is quadratic, agnostic otherwise),
\code{LSAgnostic} (the open-loop $2/(t+2)$ rule), \code{LSExact} (requires a
quadratic total objective). The closed form \eqref{eq:ls} holds when the father is
quadratic (\code{DQuadFunction} \emph{or} \code{QuadFunction}) and --- in mode
$\Ptwo$ --- for any $h_j$ (the children term is linear in $\gamma$); the father's
quadratic structure (diagonal $+$ off-diagonal) is cached in \code{set\_Block}.

\paragraph{Convergence.} The guarantee is for the \emph{smooth} case
(differentiable father). With a nondifferentiable father objective (e.g.\
\code{PolyhedralFunction}) Frank-Wolfe runs but may not reach the optimum:
\code{get\_lb}/\code{get\_ub} provide a valid bracket that may not close. This is
consistent with the literature on the nonsmooth case~\cite{white93,bolte22,%
thekum20,cheungli17}; no smoothing is implemented.

%==========================================================================%
\section{\code{Modification} handling and warm-start}\label{sec:mods}
%==========================================================================%

\code{FrankWolfeSolver} handles \emph{lazily} the \code{Modification} coming from
the inner Blocks, keeping the cached structures consistent and warm-starting the
active set / the iterate between one \code{compute()} and the next. The logic
mirrors that of \code{LagBFunction}, but is simpler because
\code{FrankWolfeSolver} is at the top of the chain: it \emph{consumes} the
Modification, it does not re-translate them into \code{FunctionMod} for an external
solver.

\paragraph{Mechanics.}
Between two \code{compute()} the external Modification are \emph{queued}.
\code{process\_modifications()}, at the start of \code{compute()}, drains the queue
and categorizes it. \code{inhibit\_Modification(true)} is reused as
\emph{play\_dumb} \emph{only} around the algorithm and the restore: the scatter
(and the final restore) change the children's objectives and generate
``self-inflicted'' Modification to be ignored, but the inhibit discards only the
new incoming ones, not the external queue already accumulated. At the end of
\code{compute()} the children's objectives go back to the original $c_0$
(\emph{restore-at-end}), so an external change acts on a clean state; at teardown
the restore uses \code{eNoMod} (the other \code{:Solver} might no longer exist).

\paragraph{Categorization (\code{LagBFunction}-style relax test).}
\code{guts\_of\_process\_modifications} produces a \emph{bit-mask}, and --- as in
\code{LagBFunction} --- feasibility is re-checked \emph{only} when the change can
\emph{shrink} the region:
\begin{itemize}[nosep]
  \item \textbf{father objective} $\to$ re-cache of the quadratic structure; atoms
        untouched;
  \item \textbf{child objective} $\to$ re-snapshot of $c_0$ and recomputation of
        the atoms' costs $c_i$, which are kept;
  \item \textbf{structure} (\code{C05FunctionModVars}, \code{NBModification},
        variable fixing) $\to$ full re-analysis $+$ drop of the active set;
  \item \textbf{feasible region} (\code{RowConstraintMod}, \code{VariableMod},
        \code{ConstraintMod}, \code{BlockModAD}, \code{BlockMod}) $\to$ \emph{0} if
        it relaxes, \emph{check}/reset if it restricts.
\end{itemize}
Variable \emph{fixing} is routed to a \emph{reset} (not to the feasibility check):
\code{ColVariable::is\_feasible()} checks sign/integrality but \emph{not} the
fixing (a \emph{state}, not a constraint), so a check would not drop the atoms
that violate it. As in \code{LagBFunction}'s TODO, the direction of a
\code{RowConstraintMod} cannot be deduced without the previous value: that case is
always re-checked.

\paragraph{\code{intHandleMod}.}
\code{eModReset} (default) always does the worst case (re-analysis $+$ drop of the
active set on any change, no warm-start); \code{eModFine} updates only the affected
cache and \emph{keeps the warm-start}.

\paragraph{Warm-start.} If a previous \code{compute()} left an active set / iterate
still valid, it restarts from there. Vanilla has no atoms (it cannot rebuild
$\bar c$ exactly) but still reuses the iterate $x$, reinitializing
$\bar c = C_{\mathrm{lit}}(x)$ (exact for linear children). The
\code{feasibility\_check} (eModFine, region change) writes each atom, calls
\code{is\_feasible()}, drops the infeasible ones and rebuilds $x$ from the
survivors (a convex combination of feasible points $\Rightarrow$ feasible).
\emph{Aggregate guard}: the $c_i$ of an aggregated atom is the convex combination
of the costs of the (lost) merged vertices, not $h(\cdot)$ at the point; with a
bounded active set and quadratic children a child-objective change cannot recompute
it exactly, so the warm-start is discarded.

\paragraph{Infeasibility propagation.} If an LMO returns \code{kInfeasible}, the
product region is empty and \code{compute()} returns \code{kInfeasible}.

%==========================================================================%
\section{Novel elements}\label{sec:novelty}
%==========================================================================%

\begin{enumerate}[leftmargin=1.4em]
\item \textbf{Two problems from a single decomposition (Sec.~\ref{sec:twoproblems}).}
  The same oracle, the same scheme and the same vertices serve both the ``true''
  composite objective $\Pone$ and the DW / perspective-cut bound $\Ptwo$, with
  $\Ptwo\ge\Pone$, selectable by a single parameter and at zero oracle cost.
  \code{FrankWolfeSolver} thus becomes a decomposition alternative to an explicit
  DW / P-C reformulation.
\item \textbf{Unified composite objective (Sec.~\ref{sec:composite}).} The three
  semantics \code{LMOLinear}/\code{LMOQuad}/\code{LMOFull} are $(\alpha,\beta)$ of
  a single scheme, with a single gap criterion \eqref{eq:gap} valid in all modes.
  The oracle handles exactly the child's separable part $h_j$.
\item \textbf{Lazy \code{Modification} handling with warm-start (Sec.~\ref{sec:mods}).}
  Porting of \code{LagBFunction}'s machinery (relax test, play\_dumb,
  restore-at-end) adapted to a solver at the top of the chain, with warm-start of
  both the active set and the vanilla iterate, and infeasibility propagation.
\item \textbf{``All configuration'' SMS++ integration.} The cross-check between
  \code{FrankWolfeSolver} and a monolithic \code{:MILPSolver} on the same Block is
  driven entirely by \code{Configuration} files: no assumption, in the test code,
  about which \code{:Solver} are attached.
\end{enumerate}

%==========================================================================%
\section{Experimental validation}\label{sec:valid}
%==========================================================================%

\paragraph{Cross-check $\Ptwo$ $=$ DPForm$+$P/C on \code{ThermalUnitBlock}.}
On unit-commitment instances, the quadratic father over a \code{ThermalUnitBlock}
with \code{ThermalUnitDPSolver} as LMO, in mode \code{LMOFull}$+$\code{eObjCvxComb}
($\Ptwo$), produces \emph{the same bound} (up to tolerance) as a \code{MILPSolver}
on the continuous-relaxed, cut-separated DP formulation with perspective cuts ---
including the fractional-commitment cases where $\Pone\neq\Ptwo$. This empirically
verifies the perspective-cut reading of $\Ptwo$ (Proposition~\ref{prop:jensen} and
following).

\paragraph{Randomized \code{Modification} on \code{MCFBlock}.}
The \code{test-mcf} tester (in the style of \code{tests/MCF\_MILP}) builds $k$
\code{MCFBlock} under a \code{DQuadFunction} father and, in randomized rounds,
changes the \emph{costs} (an objective change), the \emph{capacities} (a region
change) or \emph{fixes/unfixes} an arc (a \code{VariableMod} that shrinks/relaxes),
cross-checking against a \code{:MILPSolver} at each round, including infeasibility
propagation when closing an arc makes the MCF infeasible.

\paragraph{Regression suite.}
The regression (static $+$ \code{-M} rounds) runs vanilla / Away-step / BPCG, a
bounded active set, parallel LMOs, the \code{PolyhedralFunction} father (two-Block
bracket) and the region-\code{Modification} rounds on \code{MCFBlock}: all cases
pass. The matrix (variants $\times$ \code{eModReset}/\code{eModFine} $\times$
\code{DQuad}/\code{Quad}, and cost/cap/fix $\times$ seed $\times$ instances) is
described in the tester README.

%==========================================================================%
\section{Parameters}\label{sec:params}
%==========================================================================%

\begin{center}
\begin{tabular}{@{}lll@{}}
\toprule
Parameter & Values & Meaning\\
\midrule
\code{intLMOObj}    & 0/1/2 & \code{LMOLinear}/\code{LMOQuad}/\code{LMOFull} (Sec.~\ref{sec:composite})\\
\code{intAlgorithm} & 0/1/2 & vanilla / Away-step / BPCG\\
\code{intLineSearch}& 0/1/2 & auto / agnostic / exact\\
\code{intMaxAtoms}  & $\ge 0$ & active-set bound (0 $=$ unbounded)\\
\code{intCvxComb}   & 0/1 & \code{eObjAtX} ($\Pone$) / \code{eObjCvxComb} ($\Ptwo$, default)\\
\code{intHandleMod} & 0/1 & \code{eModReset} (default) / \code{eModFine}\\
\code{intMaxThread} & $\ge 1$ & parallel LMOs (default 1)\\
\bottomrule
\end{tabular}
\end{center}

\noindent Example \code{ComputeConfig} (mode $\Ptwo$, Away-step, warm-start):
\begin{lstlisting}[style=cfg]
intAlgorithm   1     # Away-step Frank-Wolfe
intLMOObj      2     # LMOFull: oracle handles c_j + g_j + q_j
intCvxComb     1     # eObjCvxComb: solves (P2) = DW / perspective-cut
intHandleMod   1     # eModFine: warm-start across consecutive compute()
\end{lstlisting}

%==========================================================================%
\section{Conclusions and future work}\label{sec:concl}
%==========================================================================%

\code{FrankWolfeSolver} brings the Frank-Wolfe family into SMS++ as a
decomposition \code{CDASolver}, integrated ``all by configuration''. The
conceptual contribution is that a single decomposition, with the same oracle,
solves \emph{two distinct problems} --- the composite objective over $\conv(X)$
and its DW / perspective-cut relaxation $\Ptwo\ge\Pone$ --- separated by a single
parameter and by the value/line-search bookkeeping alone, not by the oracle.
Future work: the exact line search of $\Pone$ with nonlinear $h_j$; smoothing for
the nonsmooth case; the full LDS-style forwarding of the per-child
\code{BlockSolverConfig}; benchmarking $\Ptwo$ against explicit DW / P-C
reformulations.

%==========================================================================%
\begin{thebibliography}{9}
%==========================================================================%
\bibitem{fwsurvey} G.~Braun, A.~Carderera, C.~W.~Combettes, H.~Hassani,
  A.~Karbasi, A.~Mokhtari, S.~Pokutta.
  \emph{Conditional Gradient Methods}. arXiv:2211.14103, 2022.
\bibitem{fwjl} M.~Besan\c{c}on, A.~Carderera, S.~Pokutta.
  \emph{FrankWolfe.jl: A High-Performance and Flexible Toolbox for Frank-Wolfe
  Algorithms and Conditional Gradients}. INFORMS J. on Computing, 2022.
\bibitem{white93} D.~J.~White. \emph{Extension of the Frank-Wolfe algorithm to
  concave nondifferentiable objective functions}. J.~Optim.~Theory Appl., 1993.
\bibitem{bolte22} J.~Bolte, C.~W.~Combettes, E.~Pauwels.
  \emph{The Iterates of the Frank-Wolfe Algorithm May Not Converge}.
  Mathematics of Operations Research, 2022.
\bibitem{thekum20} K.~K.~Thekumparampil, P.~Jain, P.~Netrapalli, S.~Oh.
  \emph{Projection Efficient Subgradient Method and Optimal Nonsmooth
  Frank-Wolfe Method}. NeurIPS, 2020.
\bibitem{cheungli17} E.~Cheung, Y.~Li. \emph{Nonsmooth Frank-Wolfe using Uniform
  Affine Approximations}. 2017.
\end{thebibliography}

\end{document}
